\hypertarget{experimental-setup}{%
\section{Experimental Setup}\label{experimental-setup}}

\label{sec:7}

This section describes the data, predictors, metrics, and protocols used to answer RQ1--RQ4 (Section~\ref{sec:8}).
The design follows one overriding principle, carried from the framework's independence guarantee
(Section~\ref{sec:5.3}): every predictor is evaluated against the
same simulator-derived ground truth produced by an independent process, so the claims we make are
\emph{comparative} --- which modeling choices perform better under identical conditions --- rather than
assertions of absolute accuracy in operational deployments.

\hypertarget{datasets}{%
\subsection{Datasets}\label{datasets}}

\label{sec:7.1}

\textbf{Synthetic suite.} We evaluate on seven synthetic pub-sub scenarios spanning distinct deployment
domains --- autonomous vehicles, high-frequency trading, clinical healthcare integration, centralized
hub-and-spoke enterprise systems, distributed IoT smart-city telemetry, cloud-native microservices,
and large-scale enterprise pub-sub. The scenarios are produced by a statistical topology generator
and range from 50 to 300 applications per scenario, exercising fan-out-dominated, dense-pub-sub, and
anti-pattern/SPOF regimes with different dominant failure mechanisms.

\textbf{Real-world open-source suite.} To test operational generalizability on authentic software graphs,
we evaluate SaG on three real-world open-source software architectures:
1. \textbf{Autoware.universe (ROS 2 Autonomous Driving Platform) \cite{kato2018autoware}:} An authentic cyber-physical ROS 2 pub-sub
architecture comprising 32 Applications (perception, sensing, localization, planning, control), 24 Topics with
explicit DDS QoS profiles (\texttt{RELIABLE}/\texttt{BEST\_EFFORT}, \texttt{TRANSIENT\_LOCAL}/\texttt{VOLATILE}), 3 Brokers (CycloneDDS, FastDDS, Zenoh),
6 Deployment Nodes, 10 Shared C++ Libraries (\texttt{autoware\_universe\_utils}, \texttt{tier4\_autoware\_utils}), and realistic SonarQube code metrics.
2. \textbf{Production Cloud-Native Microservices Mesh \cite{google2024onlineboutique}:} An authentic microservice architecture based on the
Google Online Boutique e-commerce benchmark, comprising 22 Microservices (order, payment, inventory, auth,
analytics, notifications), 20 Topics across Kafka, RabbitMQ, Redis PubSub, and NATS, 6 Kubernetes/Cloud nodes, and 8 shared helper libraries.
3. \textbf{Train-Ticket Railway Booking Mesh \cite{zhou2021fault}:} An authentic microservice architecture based on the Fudan
University Train-Ticket benchmark \cite{zhou2021fault}, comprising 41 Microservices (order, travel, preserve, route, seat,
payment, food, security, admin), 30 Topics, 3 Brokers (RabbitMQ, Redis PubSub, Spring Eureka naming
server), 8 deployment Nodes, and 8 shared Spring/MyBatis libraries. At 90 components it is the largest
of the three real-world graphs.

Pooled across all synthetic and real-world scenarios, the evaluation corpus exercises 1,770 components.
Every scenario is versioned in the replication package. The synthetic suite is registered in a
manifest carrying a canonical SHA-256 per dataset, and a regression test re-generates each synthetic
dataset from its configuration and fails on any divergence, so the corpus is byte-reproducible from
its configs. The three real-world graphs are not part of that manifest: they are versioned files
produced by a hand-written adapter rather than by the statistical topology generator, so the
byte-identity guarantee applies to the synthetic suite only.

All seven scenarios are used for the predictor evaluation (Section~\ref{sec:8.1}--Section~\ref{sec:8.3}), the analyses of Section~\ref{sec:5.4}--Section~\ref{sec:5.5}, and
the remediation evaluation of Section~\ref{sec:6.7}, the last of which became tractable at the largest scale by
parallelising counterfactual verification across candidates.

\hypertarget{predictors-and-baselines}{%
\subsection{Predictors and Baselines}\label{predictors-and-baselines}}

\label{sec:7.2}

The evaluation compares predictors spanning the interpretability--capacity spectrum, all consuming
the same structural analysis of each scenario:

\begin{longtable}[]{@{}
  >{\raggedright\arraybackslash}p{(\columnwidth - 4\tabcolsep) * \real{0.37}}
  >{\raggedright\arraybackslash}p{(\columnwidth - 4\tabcolsep) * \real{0.43}}
  >{\raggedright\arraybackslash}p{(\columnwidth - 4\tabcolsep) * \real{0.20}}@{}}
\caption{Predictors and baselines, and the factor each contrast isolates.\label{tab:14}}\tabularnewline
\toprule
Predictor & Description & Role \\
\midrule
\endfirsthead
\toprule
Predictor & Description & Role \\
\midrule
\endhead
\textbf{RMAV / \(Q\)} & deterministic multi-dimensional composite (Section~\ref{sec:4}) & interpretable predictor \\
\textbf{HGL} & heterogeneous graph transformer, QoS-masked & learned predictor (typed) \\
\textbf{HGL-QoS} & heterogeneous graph transformer, QoS-encoded & learned predictor (typed + QoS) \\
\textbf{GL / GL-QoS} & homogeneous GAT on the type-collapsed projection & learning baseline (untyped) \\
\textbf{Topo-BL / Topo-QoS} & structural centrality (betweenness, articulation points; QoS-weighted) & non-learning baseline \\
\bottomrule
\end{longtable}

The contrast \texttt{Topo-*} vs learned isolates the value of learning (RQ1); \texttt{GL} vs \texttt{HGL} isolates the
value of \emph{typed} heterogeneity; \texttt{HGL} vs \texttt{HGL-QoS} isolates the value of explicit QoS encoding
(RQ3); and \texttt{RMAV/Q} vs the learned predictors isolates when interpretable attribution suffices. The
structural baselines' features are kept decoupled from the GNN inputs so that no comparison leaks
information across the predictor boundary.

\hypertarget{evaluation-metrics}{%
\subsection{Evaluation Metrics}\label{evaluation-metrics}}

\label{sec:7.3}

We report metrics in three families, plus the stratification and significance machinery:

\begin{itemize}
\tightlist
\item
  \textbf{Ranking.} Spearman rank correlation \(\rho\) between predicted criticality and \(I^*(v)\) is the
  primary metric, complemented by NDCG@10 and Top-5/Top-10 overlap for the practically relevant case
  in which only a few components can be hardened.
\item
  \textbf{Identification.} Precision, recall, and F1 for critical-component detection, plus SPOF-F1 for
  articulation-point classification against simulated availability impact.
\item
  \textbf{Regression.} RMSE and MAE between predicted and simulated scores, for calibration.
\item
  \textbf{Stratified reporting.} Following the consistency check of Section~\ref{sec:5.5}, \(\rho\) is always reported \emph{by
  node type} in addition to (not instead of) any pooled figure.
\item
  \textbf{Statistical rigor.} Bootstrap 95\% confidence intervals \cite{efron1993introduction} (\(B = 2000\) resamples) on mean
  \(\rho\) \cite{spearman1904proof}, and paired Wilcoxon signed-rank tests \cite{wilcoxon1945individual} (\(p < 0.05\)) for predictor comparisons
  across scenarios.
\end{itemize}

\textbf{Validation gate thresholds.} The implementation carries two independent gate registries --- a fixed
nine-check one and the topology-class-adjusted five-check one below; Table~\ref{tab:15} and Section~\ref{sec:8.5} document only
the latter, which is the one the reported evaluation runs against. It is parameterised by topology
class, because a threshold that discriminates on a dense topology is either trivial or unattainable
on a sparse one. The default targets are \(\rho \ge 0.70\) and \(F1 \ge 0.80\); the \texttt{sparse} class, which
covers every graph in Section~\ref{sec:8.5}, applies the five checks below. We tabulate them because Section~\ref{sec:8.5} reports gate
\emph{failures}, and a reader cannot evaluate a failure against an unstated bar:

\begin{longtable}[]{@{}lrrrr@{}}
\caption{Topology-class validation gate thresholds. A graph's class is assigned from the density and hub ratio of its structural graph. All five checks must pass for the gate to pass.\label{tab:15}}\tabularnewline
\toprule
Check & \texttt{sparse} & \texttt{medium} & \texttt{dense} & \texttt{hub\_spoke} \\
\midrule
\endfirsthead
\toprule
Check & \texttt{sparse} & \texttt{medium} & \texttt{dense} & \texttt{hub\_spoke} \\
\midrule
\endhead
Spearman \(\rho\) vs simulated impact & \(\ge 0.75\) & \(\ge 0.80\) & \(\ge 0.82\) & \(\ge 0.85\) \\
\(F_1@K\) (critical-set overlap) & \(\ge 0.65\) & \(\ge 0.70\) & \(\ge 0.72\) & \(\ge 0.75\) \\
SPOF-F1 (articulation points vs \(I > 0.3\)) & \(\ge 0.60\) & \(\ge 0.65\) & \(\ge 0.65\) & \(\ge 0.70\) \\
FTR, false target rate (lower is better) & \(\le 0.30\) & \(\le 0.25\) & \(\le 0.25\) & \(\le 0.20\) \\
Predictive gain over degree centrality & \(\ge 0.02\) & \(\ge 0.03\) & \(\ge 0.03\) & \(\ge 0.03\) \\
\bottomrule
\end{longtable}

FTR is the fraction of components predicted highly vulnerable (\(V(v) > 0.60\)) whose simulated
reachability impact is nonetheless negligible (\(< 0.10\)) --- a precision check on the Vulnerability
dimension specifically. Note that the class adjustment is not uniformly a tightening relative to the
\(\rho \ge 0.70\) / \(F1 \ge 0.80\) defaults: \(\rho\) rises to \(0.75\) on the sparse class, whereas the
F1-family thresholds are \emph{relaxed}, because on a sparse graph a handful of components carries the
whole signal and a single disagreement moves F1 sharply. We tabulate the rule rather than describe it
as a tightening throughout.

\textbf{One evaluation contract, one sample.} Every variant in every table is scored by the same function
on the same node set. This is a correction rather than a description of prior practice: an earlier
version of this study scored the two predictor families on different populations and different
samples, and correcting it raised every learned variant by 0.35--0.48 \(\rho\) in-distribution while
leaving the baselines essentially unchanged (Section~\ref{sec:9.2}). Three properties now hold:

\begin{enumerate}
\def\labelenumi{\arabic{enumi}.}
\tightlist
\item
  \textbf{The evaluation key set is a function of the graph and the labels only} --- never of any variant's
  predictions --- so all variants in a cell see an identical sample. The node population is an
  explicit, recorded parameter (\texttt{application} by default, matching the claim that topology predicts
  \emph{application-layer} criticality).
\item
  \textbf{The reported figure is held-out.} All variants share one train/validation/test split pinned by
  node identity. A full-population score flatters a trained model by including the nodes it was
  fitted on while leaving a training-free baseline unchanged; the transductive figure is retained
  separately rather than reported as the headline.
\item
  \textbf{A variant that cannot cover the declared population fails loudly} rather than silently
  shrinking the sample to a per-variant subset.
\end{enumerate}

\textbf{Absent is not zero.} A stratum whose predictions or labels are constant has an \emph{undefined} rank
correlation and is reported as such, never as \(0.0\). This matters for coverage: Topic and physical
Node components carry no simulated ground truth at all (Section~\ref{sec:7.5}), and reporting them as \(0.0\) presented
a labelling gap as a measured model failure.

\hypertarget{protocols}{%
\subsection{Protocols}\label{protocols}}

\label{sec:7.4}

Two evaluation regimes are used, each answering a different generalization question.

\textbf{In-distribution (per-scenario).} For each scenario, predictors are computed and compared against
that scenario's simulated ground truth. This is the regime for RQ1 and RQ2 (Section~\ref{sec:8}): it asks how well
the attribution and learned predictors recover the criticality ordering of a \emph{known} system.

\textbf{Inductive (Leave-One-Scenario-Out).} To test generalization to \emph{unseen} architectures --- the true
pre-deployment condition --- we use Leave-One-Scenario-Out (LOSO) cross-validation, which closes the
transductive-leakage gap for the learned predictor. For each held-out scenario \(k\), the model
is trained on the remaining six scenarios (with the largest by \(|V|\) used for early stopping) and
evaluated on \(k\), whose nodes never participate in any forward pass and whose labels never enter any
loss. Results are aggregated as per-fold mean \(\pm\) std across seeds, then cross-fold mean \(\pm\)
std, with per-node-type \(\rho\) retained.

\textbf{Multi-seed.} Every configuration is run over five seeds \(\{42, 123, 456, 789, 2024\}\); reported
scores are seed means, and the across-seed standard deviation \(\sigma_{\text{seed}}\) is both
reported and reused as the noise scale in the remediation acceptance criterion (Section~\ref{sec:6.4}). The one
exception is the remediation sweep of Section~\ref{sec:6.7}, which uses the first three of these seeds for the compute
reason stated there.

\hypertarget{ground-truth-three-oracles-and-what-they-can-each-support}{%
\subsection{Ground Truth: Three Oracles, and What They Can Each Support}\label{ground-truth-three-oracles-and-what-they-can-each-support}}

\label{sec:7.5}

The three oracles introduced in Section~\ref{sec:5.1} are constructed differently, measure different quantities, and
are not interchangeable; conflating them is the most likely way to over-read a result in this paper.
This section fixes which analysis rests on which, quantifies how far they agree, and states the
label-coverage bounds that apply to each.

\begin{longtable}[]{@{}
  >{\raggedright\arraybackslash}p{(\columnwidth - 6\tabcolsep) * \real{0.25}}
  >{\raggedright\arraybackslash}p{(\columnwidth - 6\tabcolsep) * \real{0.25}}
  >{\raggedright\arraybackslash}p{(\columnwidth - 6\tabcolsep) * \real{0.25}}
  >{\raggedright\arraybackslash}p{(\columnwidth - 6\tabcolsep) * \real{0.25}}@{}}
\caption{The three simulation oracles, what each measures, and which results rest on which.\label{tab:16}}\tabularnewline
\toprule
Symbol & Engine & Quantity & Used for \\
\midrule
\endfirsthead
\toprule
Symbol & Engine & Quantity & Used for \\
\midrule
\endhead
\(I^*(v)\) & \texttt{FaultInjector} & Mean subscriber feed-loss fraction under a BFS cascade & Learned-predictor labels; Tables~\ref{tab:18} and~\ref{tab:20} (Section~\ref{sec:8.1}); the sensitivity sweeps of Section~\ref{sec:8.3} \\
\(I_{\text{comp}}(v)\) & \texttt{FailureSimulator} & \(0.35\,\text{reachability} + 0.25\,\text{fragmentation} + 0.25\,\text{throughput} + 0.15\,\text{flow}\) & Validation gates; the RMAV dimension decomposition; Section~\ref{sec:5.4} and Section~\ref{sec:5.5}; remediation acceptance (Section~\ref{sec:6.4}) \\
\(I_{\text{dyn}}(v)\) & \texttt{MessageFlowSimulator} & Delivery-rate loss suffered by \emph{surviving} consumers, by discrete-event simulation of traffic & Reported construct-validity check only --- no labels, no gates, no tables \\
\bottomrule
\end{longtable}

The two cascade oracles run with a step-function blast-semantics propagation scheme (probability
\(1.0\) for library cascade), \texttt{propagation\_threshold} default \(0.2\), a \(10\)-epoch horizon, and the
five seeds of Section~\ref{sec:7.4}, \(\{42, 123, 456, 789, 2024\}\). \(I_{\text{dyn}}\) shares the seed set and runs
\(60\) simulated seconds per component, with the fault injected at the midpoint.

\textbf{Measured agreement between the two cascade oracles is weak.} Their scales differ, so only rank
agreement is meaningful; across the seven scenarios, mean Spearman \(\rho = 0.394\) and mean top-20\%
Jaccard \(= 0.286\), ranging from \(\rho = 0.578\) (Enterprise) down to \(\rho = 0.092\) (Hub-and-Spoke,
where they are effectively uncorrelated). All seven correlations are positive, which is a weak
convergent-validity argument --- two differently-constructed simulators do agree directionally, so
neither is purely an artifact of its own construction --- but at \(\rho \approx 0.39\) it is weak, and we
apply the resulting constraint to our own analyses: a result established against one oracle is not
evidence for a claim measured against the other. Section~\ref{sec:8.2} flags where this bites.

\textbf{\(I_{\text{dyn}}\) agrees with \(I^*\) far more strongly, and --- crucially --- does not share its worst
case.} Mean Spearman \(\rho(I_{\text{dyn}}, I^*) = 0.765\), minimum \(0.548\) (Hub-and-Spoke) --- against
mean \(0.394\), minimum \(0.092\) for the two topological oracles above. Hub-and-Spoke is precisely where
\(I^*\) and \(I_{\text{comp}}\) collapse to near-independence; \(I_{\text{dyn}}\) still agrees with \(I^*\)
there at \(\rho = 0.548\), its lowest agreement in the cohort but far from uncorrelated. Because
\(I_{\text{dyn}}\) reaches this ranking by simulating traffic through queues rather than by traversing
\texttt{DEPENDS\_ON}, the result is evidence of a different kind than Section~\ref{sec:7.5}'s first finding: it rules out the
cascade \emph{algorithm} as the source of \(I^*\)'s ranking, which the \(I_{\text{comp}}\) comparison alone
cannot do (Section~\ref{sec:9.2}). Top-\(K\) membership is the weaker half of this result --- mean top-20\% Jaccard is
\(0.316\), comparable to the \(0.286\) of the two topological oracles --- so this is corroboration of
\emph{ranking}, not of critical-set identification; no \(F_1@K\) claim in Section~\ref{sec:8.1} is supported by
\(I_{\text{dyn}}\).

\textbf{Label coverage and the noise ceiling.} Three further properties bound what any reported figure can
mean. First, the cascade model has no rule expressing the failure of a Topic or of a physical Node,
so those types carry no ground truth at all --- 30--47\% of components per scenario are unlabelled, they
are excluded from scoring rather than scored as zero, and predictions for them are never validated.
Broker labels are degenerate in three of seven scenarios for a related reason. Second, the three
oracles do not cover the same components, so every agreement figure above is computed over the
intersection rather than over the scenario. \(I_{\text{dyn}}\) observes only what carries pub-sub
traffic: it scores Applications and those Libraries that publish or subscribe in their own right,
and records Brokers, physical Nodes, Topics, and purely-consumed Libraries as unmeasured rather than
as harmless. On \texttt{enterprise\_system} that is 349 components against \(I^*\)'s 360 --- it gives up the ten
Brokers and one non-publishing Library, and gains nothing \(I^*\) lacks. Third, the labels
have a reproducibility ceiling: across seeds, the ground truth agrees with \emph{itself} at test--retest
\(\rho\) of 0.807--1.000, and its own top-20\% critical set agrees at Jaccard 0.44--1.00 (deterministic:
\texttt{FaultInjector}`s cascade previously iterated an unordered subscriber set while consuming seeded
random draws, so re-running the \emph{same} seed in a different process could still change the label ---
this has been fixed and is disclosed as an instrument defect in Section~\ref{sec:9.2}, and the figures here are the
post-fix, process-independent ones). No method can exceed the former, and every top-\(K\) metric
inherits the latter --- a reported \(F_1@K\) on \texttt{microservices\_system}, where the labels' own set
stability is 0.44, should not be read to a precision the labels do not have.

\hypertarget{model-configuration-and-implementation}{%
\subsection{Model Configuration and Implementation}\label{model-configuration-and-implementation}}

\label{sec:7.6}

The learned predictors are implemented in PyTorch Geometric \cite{fey2019fast}. Table~\ref{tab:17} fixes every hyperparameter; all
variants share it, so the contrasts of Section~\ref{sec:7.2} isolate architecture and features rather than tuning
budget. The values were fixed before the reported runs and not tuned per scenario --- there is no
per-scenario hyperparameter search anywhere in this study, which is a deliberate constraint (a search
per scenario would leak held-out information under the in-distribution protocol) and also a
limitation, since a tuned baseline might close some of the margins in Section~\ref{sec:8.1}.

\begin{longtable}[]{@{}
  >{\raggedright\arraybackslash}p{(\columnwidth - 2\tabcolsep) * \real{0.50}}
  >{\raggedright\arraybackslash}p{(\columnwidth - 2\tabcolsep) * \real{0.50}}@{}}
\caption{Learned-predictor configuration. Identical across HGL, \(HGL\text{-}QoS\), GL and GL-QoS except where the architecture differs by construction.\label{tab:17}}\tabularnewline
\toprule
Component & Setting \\
\midrule
\endfirsthead
\toprule
Component & Setting \\
\midrule
\endhead
Convolution & \texttt{HGTConv} (heterogeneous); \texttt{GATConv} for the homogeneous GL variants \\
Layers & 3 \\
Hidden channels & 64 \\
Attention heads & 4 \\
Dropout & 0.2 \\
Input projection & per-node-type linear → LayerNorm → ReLU \\
Output heads & four RMAV residual MLPs + one composite head, sigmoid-activated \\
Optimizer & AdamW, learning rate \(3\times10^{-4}\), weight decay \(1\times10^{-4}\) \\
LR schedule & \texttt{CosineAnnealingWarmRestarts}, \(T_0 = \max(50, \text{epochs}/4)\), \(T_{\text{mult}} = 2\), \(\eta_{\min} = 0.01\cdot\text{lr}\) \\
Gradient clipping & max-norm 1.0 \\
Epochs / early stopping & 300, patience 30 on validation loss \\
Node splits & 60\% train / 20\% validation / 20\% test, pinned by node identity (Section~\ref{sec:7.3}) \\
Loss & composite MSE \(+\ 0.5\,\)multitask \(+\ 0.3\,\)ListMLE ranking \(+\ 0.1\,\)pairwise margin \(+\ 0.1\,\)RMAV consistency \\
\bottomrule
\end{longtable}

\textbf{Hardware and runtime.} Training and evaluation were run on a single workstation; the LOSO sweep is
the dominant cost, at roughly 31 minutes for HGL and 36 for \(HGL\text{-}QoS\) across all folds and
seeds, against 5--6 minutes for the homogeneous variants and well under a minute for the training-free
baselines. The CI/CD gate measurements of Section~\ref{sec:8.4} were taken on the same machine rather than on hosted
runner hardware, so they should be read as an order-of-magnitude feasibility result rather than as a
calibrated figure for any particular CI provider.

\begin{center}\rule{0.5\linewidth}{0.5pt}\end{center}
